\section{Horizons, predictability and deadlines}
\label{sec:horizons}

Forecasts for renewable energy span nine orders of magnitude in lead time, from seconds for plant controllers to decades for financing (Figure~\ref{fig:horizons}). The industry's workhorse is the day-ahead forecast in quarter-hours, followed by updates through the delivery day and outlooks of up to two weeks for planning. Table~\ref{tab:horizons} lists the inputs and methods that dominate each horizon.

% Figure 2: horizon map. Conceptual; the ranges are listed in data/literature/horizons.csv.
\begin{figure*}[t]
\centering
\begin{tikzpicture}
\begin{axis}[paper, width=0.83\textwidth, height=6.6cm,
  xmode=log, xmin=1.5e-4, xmax=6e5,
  ymin=-0.3, ymax=7.6, ytick={1,...,7}, y dir=reverse,
  yticklabels={Plant control and hybrid storage, Grid control rooms and intraday trading, Day-ahead trading and scheduling, {O\&M, construction, commissioning}, {Hedging, hydro and fuel planning}, Budgets and re-forecasts, {Development, financing, PPAs}},
  y tick label style={font=\scriptsize,text=ewfInk},
  xtick={2.7778e-4,1.6667e-2,1,24,168,730,8760,87600},
  xticklabels={1\,s,1\,min,1\,h,1\,day,1\,week,1\,month,1\,year,10\,years},
  xlabel={Lead time (log scale)}, grid=none, xmajorgrids, ymajorgrids=false,
  axis y line*=left, clip=false]
% method regimes in the background
\fill[ewfPanel] (axis cs:1.5e-4,-0.3) rectangle (axis cs:6,7.6);
\fill[white]    (axis cs:6,-0.3) rectangle (axis cs:360,7.6);
\fill[ewfPanel] (axis cs:360,-0.3) rectangle (axis cs:5040,7.6);
\fill[white]    (axis cs:5040,-0.3) rectangle (axis cs:6e5,7.6);
\node[font=\scriptsize,text=ewfInkSecondary,anchor=north,align=center] at (axis cs:0.04,-0.3) {measurements, satellite,\\rapid-update NWP};
\node[font=\scriptsize,text=ewfInkSecondary,anchor=north,align=center] at (axis cs:46,-0.3) {NWP and\\ensembles};
\node[font=\scriptsize,text=ewfInkSecondary,anchor=north,align=center] at (axis cs:1350,-0.3) {ensemble\\anomalies};
\node[font=\scriptsize,text=ewfInkSecondary,anchor=north,align=center] at (axis cs:55000,-0.3) {historical records,\\climate data};
% one bar per user group (one hue: the rows are nominal categories)
\foreach \lo/\hi/\y in {2.7778e-4/0.5/1, 0.0833/24/2, 12/48/3, 24/240/4, 336/5040/5, 720/8760/6, 8760/262800/7}{
  \edef\tmp{\noexpand\fill[ewfC1,rounded corners=1.5pt] (axis cs:\lo,\y-0.26) rectangle (axis cs:\hi,\y+0.26);}
  \tmp
}
\end{axis}
\end{tikzpicture}
\caption{\textbf{Who forecasts how far ahead.} Each bar spans the lead times a user group typically works with; shaded bands show which information dominates. Performance analysis, lost-energy calculations and insurance claims use the past and are not shown.}
\label{fig:horizons}
\end{figure*}

% Table 3: forecast horizons.
\begin{table*}[t]
\centering\small
\caption{\textbf{Forecast horizons, their inputs and their users.}}
\label{tab:horizons}
\setlength{\tabcolsep}{4pt}
\begin{tabularx}{\textwidth}{@{}L{2.3cm}L{2.0cm}L{1.6cm}YL{3.4cm}L{3.4cm}@{}}
\toprule
Horizon & Lead time & Resolution & Main inputs & Typical methods & Main users \\
\midrule
Very short term & seconds to 30\,min & 1\,s--1\,min & sky cameras, live plant data, nacelle or lidar wind & persistence, image-based cloud tracking & plant controllers, hybrid PV and storage \\
Nowcasting & 30\,min to 6\,h & 5--15\,min & satellite images, live SCADA, rapid-update NWP, radar & cloud motion vectors, machine learning on images, blending & intraday traders, grid control rooms \\
Intraday & 1 to 24\,h & 15\,min & latest NWP runs, nowcasts, live data & statistical blending, frequent re-runs & intraday auctions and continuous trading \\
Day-ahead & 12 to 48\,h & 15\,min & 00 and 06\,UTC runs of several models & physical model plus statistical correction, quantiles & day-ahead auction, scheduling, O\&M planning \\
Medium range & 2 to 15 days & 1\,h--1 day & global NWP and ensembles & calibrated ensemble statistics & maintenance, construction, weekly hedging \\
Sub-seasonal & 2 to 6 weeks & weekly & extended-range ensembles & weekly anomalies, regime probabilities & monthly forwards, hydro, gas storage \\
Seasonal & 1 to 7 months & monthly & multi-system seasonal forecasts & tercile probabilities & budgets, reservoirs, winter adequacy \\
Long term & 1 to 30+ years & hourly to annual & reanalysis, satellite records, on-site data, climate projections & long-term correction, P-values, Monte Carlo & development, financing, PPAs \\
\bottomrule
\end{tabularx}
\end{table*}


\para{How skill decays.} Small errors in the initial state grow until the forecast is no better than climatology~\cite{lorenz1963deterministic}. For the surface variables that matter for energy, the decay is faster than for the upper-air fields used in headline scores: cloud cover and irradiance lose much of their skill within a few days, 10\,m and hub-height wind somewhat later. The case study measures this directly: for daytime 10\,m wind over Germany in summer 2025, single models kept positive skill against climatology for \valWindGfsSkillZeroDays--\valWindAifsSkillZeroDays{} days, and a bias-corrected multi-model mean for \valWindMmmBcSkillZeroDays{} days (Section~\ref{sec:case}). Beyond the deterministic limit only ensembles carry information, first as probabilities of weather regimes, at sub-seasonal range as weekly anomalies, and at seasonal range as shifted odds, which in Europe are informative mainly in winter.

\para{How accurate is good.} Published accuracy depends strongly on normalisation (installed capacity, mean production, daylight hours only) and must be compared with care. Representative German wind farms reached a day-ahead RMSE of 10--15\,\% of installed capacity, control areas 6--7\,\% and Germany as a whole 5--6\,\% in 2007~\cite{ernst2007predicting}; nationwide errors were about 3--4\,\% a decade later~\cite{ubimet2016windsage}. Aggregating 152 PV systems reduced day-ahead MAE by at least 18\,\% and RMSE by at least 20\,\% compared with forecasting each system~\cite{visser2022aggregated}. Both effects, aggregation and steady model improvement, are quantified for the case study in Section~\ref{sec:case}.

\para{Deadlines.} A forecast is only useful if it arrives before the decision. For the European day-ahead auction, which closes at 12:00\,CET on the day before delivery, the freshest global runs are those started at 00\,UTC, available in the early morning; regional models that run every three hours, and rapid-update runs such as DWD's hourly ICON-D2-RUC~\cite{wetterdienst2025ruc}, add later information if the pipeline is fast. The 12\,UTC runs serve the intraday auctions and continuous trading (Figure~\ref{fig:market}). Since 1 October 2025 the day-ahead market clears in 15-minute units~\cite{apg2025mtu}, so hourly forecasts no longer match the product.
